
Reinforcement learning from human feedback (RLHF) produces AI outputs with characteristic stylistic features, including structured lists, safety prefaces, chain-of-thought scaffolding markers, and hedging language. Whether annotators who interact repeatedly with deployed RLHF-optimized systems develop stronger preferences for these features — compared to annotators without such exposure — is an empirical question that existing preference datasets have not been used to address. This paper introduces the AI-Idiomatic Formatting Signal (AIFS) construct, a regex-based composite measure across four feature classes, and evaluates it on 80,342 preference pairs from the Anthropic HH-RLHF corpus using a conditional logit model with 27,034 semantic cluster fixed effects and a response-length covariate. AIFS features predict annotator preference (β₁ = 0.0246, p < 0.001). The adaptation interaction term — whether annotators in the deployed-AI condition prefer AIFS features more strongly than annotators in the base condition — is not statistically significant (β₄ = −0.0016, OR = 0.9984, 95\% CI [0.9861, 1.0108], p = 0.796). Examination of the dataset reveals that the HH-RLHF split labels encode data collection provenance rather than annotator exposure history, rendering them invalid proxies for the intended adaptation contrast. The null result on β₄ is therefore uninterpretable as direct evidence regarding the bidirectional alignment hypothesis. This finding identifies a structural limitation in existing RLHF preference datasets for bidirectional alignment research and motivates the collection of annotator-linked, longitudinally designed preference corpora.

